% FreqCone: Frequency Entailment Cones for Fine-Grained Generalized Category Discovery
% Working draft. Every number in Section 5 is generated from the JSON results by make_tables.py
% (macros.tex, tab_*.tex, lambda_sweep.dat). Section 6 is a protocol: nothing in it has been run.
\documentclass[10pt]{article}
\usepackage[margin=1in]{geometry}
\usepackage{amsmath,amssymb,amsthm,booktabs,graphicx,xcolor}
\usepackage[hidelinks]{hyperref}
\usepackage{algorithm,algpseudocode}
\usepackage{pgfplots}
\usepgfplotslibrary{groupplots}
\pgfplotsset{compat=1.16}
\newtheorem{proposition}{Proposition}
\newtheorem{lemma}[proposition]{Lemma}
\DeclareMathOperator{\artanh}{artanh}
\DeclareMathOperator{\sg}{sg}
\newcommand{\Lb}{\mathcal{L}}
\IfFileExists{macros.tex}{% generated by make_tables.py -- do not edit by hand
\newcommand{\selLambda}{0.03}
\newcommand{\devLama}{0}
\newcommand{\devAlla}{72.8}
\newcommand{\devOlda}{65.4}
\newcommand{\devNewa}{76.5}
\newcommand{\devPara}{43.2}
\newcommand{\devNtpa}{64.2}
\newcommand{\devValca}{89.8}
\newcommand{\devLamb}{0.03}
\newcommand{\devAllb}{74.3}
\newcommand{\devOldb}{64.4}
\newcommand{\devNewb}{79.2}
\newcommand{\devParb}{37.7}
\newcommand{\devNtpb}{66.8}
\newcommand{\devValcb}{91.7}
\newcommand{\devLamc}{0.1}
\newcommand{\devAllc}{77.7}
\newcommand{\devOldc}{75.0}
\newcommand{\devNewc}{79.1}
\newcommand{\devParc}{72.2}
\newcommand{\devNtpc}{76.1}
\newcommand{\devValcc}{80.8}
\newcommand{\devLamd}{0.3}
\newcommand{\devAlld}{63.9}
\newcommand{\devOldd}{63.6}
\newcommand{\devNewd}{64.1}
\newcommand{\devPard}{85.6}
\newcommand{\devNtpd}{91.0}
\newcommand{\devValcd}{72.7}
\newcommand{\devLame}{1}
\newcommand{\devAlle}{48.9}
\newcommand{\devOlde}{42.8}
\newcommand{\devNewe}{51.9}
\newcommand{\devPare}{100.0}
\newcommand{\devNtpe}{97.0}
\newcommand{\devValce}{55.6}
\newcommand{\hoAllGcd}{76.0}
\newcommand{\hosdAllGcd}{11.5}
\newcommand{\hoOldGcd}{78.7}
\newcommand{\hosdOldGcd}{13.7}
\newcommand{\hoNewGcd}{74.7}
\newcommand{\hosdNewGcd}{11.4}
\newcommand{\hoParGcd}{24.6}
\newcommand{\hosdParGcd}{2.1}
\newcommand{\hoNtpGcd}{22.2}
\newcommand{\hosdNtpGcd}{4.6}
\newcommand{\hoValcGcd}{85.8}
\newcommand{\hosdValcGcd}{7.4}
\newcommand{\hoAllHyp}{75.4}
\newcommand{\hosdAllHyp}{9.6}
\newcommand{\hoOldHyp}{73.1}
\newcommand{\hosdOldHyp}{14.4}
\newcommand{\hoNewHyp}{76.6}
\newcommand{\hosdNewHyp}{7.3}
\newcommand{\hoParHyp}{30.8}
\newcommand{\hosdParHyp}{12.4}
\newcommand{\hoNtpHyp}{32.2}
\newcommand{\hosdNtpHyp}{26.7}
\newcommand{\hoValcHyp}{84.2}
\newcommand{\hosdValcHyp}{3.4}
\newcommand{\hoAllFour}{70.9}
\newcommand{\hosdAllFour}{2.6}
\newcommand{\hoOldFour}{73.7}
\newcommand{\hosdOldFour}{3.3}
\newcommand{\hoNewFour}{69.5}
\newcommand{\hosdNewFour}{2.5}
\newcommand{\hoParFour}{26.0}
\newcommand{\hosdParFour}{6.2}
\newcommand{\hoNtpFour}{28.7}
\newcommand{\hosdNtpFour}{22.2}
\newcommand{\hoValcFour}{81.0}
\newcommand{\hosdValcFour}{5.0}
\newcommand{\hoAllNaive}{44.5}
\newcommand{\hosdAllNaive}{20.2}
\newcommand{\hoOldNaive}{45.9}
\newcommand{\hosdOldNaive}{25.4}
\newcommand{\hoNewNaive}{43.8}
\newcommand{\hosdNewNaive}{17.7}
\newcommand{\hoParNaive}{74.7}
\newcommand{\hosdParNaive}{31.8}
\newcommand{\hoNtpNaive}{85.1}
\newcommand{\hosdNtpNaive}{20.1}
\newcommand{\hoValcNaive}{60.8}
\newcommand{\hosdValcNaive}{17.8}
\newcommand{\hoAllWzero}{75.8}
\newcommand{\hosdAllWzero}{12.0}
\newcommand{\hoOldWzero}{72.4}
\newcommand{\hosdOldWzero}{11.1}
\newcommand{\hoNewWzero}{77.5}
\newcommand{\hosdNewWzero}{14.0}
\newcommand{\hoParWzero}{30.1}
\newcommand{\hosdParWzero}{13.0}
\newcommand{\hoNtpWzero}{28.6}
\newcommand{\hosdNtpWzero}{31.8}
\newcommand{\hoValcWzero}{82.1}
\newcommand{\hosdValcWzero}{5.8}
\newcommand{\hoAllSel}{76.5}
\newcommand{\hosdAllSel}{11.6}
\newcommand{\hoOldSel}{75.8}
\newcommand{\hosdOldSel}{6.7}
\newcommand{\hoNewSel}{76.8}
\newcommand{\hosdNewSel}{14.1}
\newcommand{\hoParSel}{29.9}
\newcommand{\hosdParSel}{13.1}
\newcommand{\hoNtpSel}{29.4}
\newcommand{\hosdNtpSel}{33.0}
\newcommand{\hoValcSel}{84.4}
\newcommand{\hosdValcSel}{7.8}
\newcommand{\hoAllWone}{69.8}
\newcommand{\hosdAllWone}{26.2}
\newcommand{\hoOldWone}{67.2}
\newcommand{\hosdOldWone}{21.3}
\newcommand{\hoNewWone}{71.2}
\newcommand{\hosdNewWone}{28.7}
\newcommand{\hoParWone}{44.5}
\newcommand{\hosdParWone}{38.2}
\newcommand{\hoNtpWone}{41.1}
\newcommand{\hosdNtpWone}{49.8}
\newcommand{\hoValcWone}{73.3}
\newcommand{\hosdValcWone}{17.3}
\newcommand{\dAllGcd}{+0.6}
\newcommand{\wAllGcd}{1/3}
\newcommand{\dNewGcd}{-1.9}
\newcommand{\wNewGcd}{1/3}
\newcommand{\dParGcd}{-6.2}
\newcommand{\wParGcd}{2/3}
\newcommand{\dNtpGcd}{-10.0}
\newcommand{\wNtpGcd}{2/3}
\newcommand{\dAllFour}{-4.5}
\newcommand{\wAllFour}{1/3}
\newcommand{\dNewFour}{-7.0}
\newcommand{\wNewFour}{0/3}
\newcommand{\dParFour}{-4.8}
\newcommand{\wParFour}{0/3}
\newcommand{\dNtpFour}{-3.6}
\newcommand{\wNtpFour}{1/3}
\newcommand{\dAllNaive}{-30.9}
\newcommand{\wAllNaive}{0/3}
\newcommand{\dNewNaive}{-32.8}
\newcommand{\wNewNaive}{0/3}
\newcommand{\dParNaive}{+43.9}
\newcommand{\wParNaive}{3/3}
\newcommand{\dNtpNaive}{+52.9}
\newcommand{\wNtpNaive}{3/3}
\newcommand{\dAllWzero}{+0.4}
\newcommand{\wAllWzero}{1/3}
\newcommand{\dNewWzero}{+0.9}
\newcommand{\wNewWzero}{1/3}
\newcommand{\dParWzero}{-0.7}
\newcommand{\wParWzero}{2/3}
\newcommand{\dNtpWzero}{-3.6}
\newcommand{\wNtpWzero}{1/3}
\newcommand{\dAllSel}{+1.1}
\newcommand{\wAllSel}{2/3}
\newcommand{\dNewSel}{+0.3}
\newcommand{\wNewSel}{2/3}
\newcommand{\dParSel}{-0.9}
\newcommand{\wParSel}{1/3}
\newcommand{\dNtpSel}{-2.8}
\newcommand{\wNtpSel}{1/3}
\newcommand{\dAllWone}{-5.6}
\newcommand{\wAllWone}{2/3}
\newcommand{\dNewWone}{-5.4}
\newcommand{\wNewWone}{2/3}
\newcommand{\dParWone}{+13.7}
\newcommand{\wParWone}{2/3}
\newcommand{\dNtpWone}{+8.9}
\newcommand{\wNtpWone}{1/3}
\newcommand{\pairSelVsWzeroAll}{+0.7}
\newcommand{\pairwSelVsWzeroAll}{1/3}
\newcommand{\pairSelVsFourAll}{+5.6}
\newcommand{\pairwSelVsFourAll}{2/3}
\newcommand{\pairSelVsFourNew}{+7.3}
\newcommand{\pairwSelVsFourNew}{2/3}
\newcommand{\pairWoneVsNaiveAll}{+25.4}
\newcommand{\pairwWoneVsNaiveAll}{3/3}
\newcommand{\pairWoneVsNaiveNtp}{-44.0}
\newcommand{\pairwWoneVsNaiveNtp}{0/3}
\newcommand{\vdAllVtwo}{+2.5}
\newcommand{\vwAllVtwo}{2/3}
\newcommand{\vdNewVtwo}{+2.8}
\newcommand{\vwNewVtwo}{2/3}
\newcommand{\vdNtpVtwo}{-19.8}
\newcommand{\vwNtpVtwo}{0/3}
\newcommand{\vdAllVthree}{+0.4}
\newcommand{\vwAllVthree}{1/3}
\newcommand{\vdNewVthree}{-1.2}
\newcommand{\vwNewVthree}{1/3}
\newcommand{\vdNtpVthree}{-18.1}
\newcommand{\vwNtpVthree}{0/3}
\newcommand{\vdAllVthreeWone}{-31.4}
\newcommand{\vwAllVthreeWone}{0/3}
\newcommand{\vdNewVthreeWone}{-31.0}
\newcommand{\vwNewVthreeWone}{0/3}
\newcommand{\vdNtpVthreeWone}{+57.7}
\newcommand{\vwNtpVthreeWone}{3/3}
\newcommand{\VthreeConeRange}{75--83\%}
\newcommand{\vdAllVthreeVsVtwo}{-2.1}
\newcommand{\vwAllVthreeVsVtwo}{1/3}
\newcommand{\vdNewVthreeVsVtwo}{-4.1}
\newcommand{\vwNewVthreeVsVtwo}{1/3}
\newcommand{\vdNtpVthreeVsVtwo}{+1.8}
\newcommand{\vwNtpVthreeVsVtwo}{3/3}
\newcommand{\sensAllWzero}{75.8}
\newcommand{\sensOldWzero}{72.4}
\newcommand{\sensNewWzero}{77.5}
\newcommand{\sensParWzero}{30.1}
\newcommand{\sensNtpWzero}{28.6}
\newcommand{\sensValcWzero}{82.1}
\newcommand{\sensdAllWzero}{+0.4}
\newcommand{\senswAllWzero}{1/3}
\newcommand{\sensdNewWzero}{+0.9}
\newcommand{\senswNewWzero}{1/3}
\newcommand{\sensdNtpWzero}{-3.6}
\newcommand{\senswNtpWzero}{1/3}
\newcommand{\sensAllWthreec}{76.5}
\newcommand{\sensOldWthreec}{75.8}
\newcommand{\sensNewWthreec}{76.8}
\newcommand{\sensParWthreec}{29.9}
\newcommand{\sensNtpWthreec}{29.4}
\newcommand{\sensValcWthreec}{84.4}
\newcommand{\sensdAllWthreec}{+1.1}
\newcommand{\senswAllWthreec}{2/3}
\newcommand{\sensdNewWthreec}{+0.3}
\newcommand{\senswNewWthreec}{2/3}
\newcommand{\sensdNtpWthreec}{-2.8}
\newcommand{\senswNtpWthreec}{1/3}
\newcommand{\sensAllWone}{69.8}
\newcommand{\sensOldWone}{67.2}
\newcommand{\sensNewWone}{71.2}
\newcommand{\sensParWone}{44.5}
\newcommand{\sensNtpWone}{41.1}
\newcommand{\sensValcWone}{73.3}
\newcommand{\sensdAllWone}{-5.6}
\newcommand{\senswAllWone}{2/3}
\newcommand{\sensdNewWone}{-5.4}
\newcommand{\senswNewWone}{2/3}
\newcommand{\sensdNtpWone}{+8.9}
\newcommand{\senswNtpWone}{1/3}
\newcommand{\logAllHyp}{73.6}
\newcommand{\logParHyp}{34.6}
\newcommand{\logNtpHyp}{54.0}
\newcommand{\logDepthHyp}{4.60}
\newcommand{\logAllVzero}{56.7}
\newcommand{\logParVzero}{78.8}
\newcommand{\logNtpVzero}{88.6}
\newcommand{\logDepthVzero}{1.25}
\newcommand{\logAllVone}{29.9}
\newcommand{\logParVone}{97.7}
\newcommand{\logNtpVone}{96.5}
\newcommand{\logDepthVone}{0.89}
\newcommand{\logAllVtwo}{48.9}
\newcommand{\logParVtwo}{100.0}
\newcommand{\logNtpVtwo}{97.0}
\newcommand{\logDepthVtwo}{4.60}
}{}
% placeholders for result prose, filled in after the sweep

\title{FreqCone: Frequency Entailment Cones for\\ Fine-Grained Generalized Category Discovery}
\author{Anonymous}
\date{}

\begin{document}
\maketitle

\begin{abstract}
Generalized Category Discovery (GCD) must recognise known classes and discover novel ones in unlabelled data.
Strong fine-grained GCD methods exploit class hierarchy, taken from external semantics, from noisy hierarchical
pseudo-labels, or left implicit in hyperbolic geometry. We observe that every image already carries a
deterministic, label-free hierarchy: low-pass filtering at increasing cut-off frequencies merges images that
differ only in fine detail, which yields nested partitions of the data. FreqCone turns this frequency
filtration into hyperbolic entailment. The low-passed view of an image must contain the image in its
entailment cone, and the blur level is encoded as depth in the Poincar\'e ball rather than as a slice of
latent dimensions. We prove that the filtration is a hierarchy, that slicing latent dimensions after a
full-vector exponential map entangles depth, that zero cone energy implies depth ordering, and that cones bound
the angle between an image and its blurred ancestors. A controlled mechanism study with 48 classes in two
levels, one development seed and three held-out seeds, shows that the entailment weight trades fine-grained
clustering accuracy for hierarchical consistency, and that a weight chosen on known-class validation data
matches HypCD within seed variance (\dAllSel{} points of All accuracy over three held-out seeds). The cone acts
only when sharp views lie deeper than their blurred ancestors, as the analysis predicts; otherwise it is vacuous
or pushes samples away from their ancestors. On fresh seeds, letting the cone move the sharp-view norm activates
it, at the price of fine-grained accuracy. In this setting FreqCone is a controllable trade-off rather than an
accuracy gain.
We release a drop-in PyTorch module and a protocol, with decision criteria fixed in advance, for the
Semantic Shift Benchmark.
\end{abstract}

% =====================================================================================================
\section{Introduction}
Generalized Category Discovery (GCD)~\cite{vaze2022gcd} gives a model a labelled set that covers some of the
classes and an unlabelled set that contains those classes and new ones; the model must group the unlabelled
images into all classes. The fine-grained splits of the Semantic Shift Benchmark (SSB)~\cite{vaze2022ssb}
(CUB, Stanford Cars, FGVC-Aircraft) are hard because a novel class often differs from a known one only in
small details while sharing its coarse appearance: the body shape of a bird family, the silhouette of a car
make, the airframe of an aircraft family.

Recent methods exploit this shared coarse structure. SelEx~\cite{rastegar2024selex} builds hierarchical
pseudo-labels by clustering at several granularities; SEAL~\cite{he2025seal} uses the semantic hierarchy that
comes with each dataset; language-based methods such as ATCG~\cite{atcg2026} borrow concepts from text.
HypCD~\cite{liu2025hypcd} maps features into the Poincar\'e ball, where tree-like structure embeds with low
distortion, but leaves the hierarchy implicit. FourEx~\cite{rastegar2026fourex} supervises with
Fourier-filtered views and assigns frequency bands to slices of the latent vector.

We ask whether the image itself can supply a hierarchy that is label-free, deterministic rather than
estimated, and compatible with hyperbolic geometry. Low-pass filtering satisfies the first two requirements.
Blurring removes the details that separate sibling classes before it removes the layout shared by their
parent, and the filters at increasing cut-offs form a filtration whose induced partitions are nested
(Proposition~\ref{prop:filtration}). For the third requirement we show that latent slicing interacts badly
with the full-vector exponential map of hyperbolic GCD (Proposition~\ref{prop:slice}). We therefore encode the
blur level as depth: the blurred view is an ancestor whose entailment cone~\cite{ganea2018cones} must contain
the sharp view. Khrulkov et al.~\cite{khrulkov2020hyperbolic} observed that ambiguous images tend to lie near
the origin of a learned Poincar\'e embedding; FreqCone makes a related ordering an explicit training signal.

A second motivation is practical. In retail shelf monitoring, a new product variant should be reported as
``unknown item of known family X'' rather than only ``unknown''. Because coarse views of siblings share a
direction, a novel sample can be attributed to the family of the nearest known class
(Section~\ref{sec:attr}); we measure this as novel-to-parent attribution.

\paragraph{Contributions.}
\begin{enumerate}
\item FreqCone, a training objective that turns the low-pass filtration of each image into a chain of
hyperbolic entailment constraints. Coarse views sit at fixed depths and two stop-gradients prevent the
collapse modes we observed during development (Section~\ref{sec:method}). The objective adds two forward
passes on blurred views during training and nothing at inference.
\item An analysis: the filtration is a hierarchy; slicing after the exponential map entangles depth; zero cone
energy implies depth ordering; cones bound the angle between a sample and its ancestors
(Section~\ref{sec:analysis}). Every statement is also checked numerically in the released code.
\item A controlled mechanism study that reports every variant we tried, including the failed ones, measures
the trade-off set by the entailment weight, and identifies the depth condition under which the cone acts
(Section~\ref{sec:mech}); and a protocol with decision criteria for the SSB evaluation (Section~\ref{sec:ssb}).
\end{enumerate}

% =====================================================================================================
\section{Related work}
\paragraph{Generalized category discovery.} GCD~\cite{vaze2022gcd} combines supervised and self-supervised
contrastive learning with semi-supervised $k$-means. SimGCD~\cite{wen2023simgcd} trains a parametric
classifier with self-distillation; $\mu$GCD~\cite{vaze2023mugcd} uses a mean-teacher. Recent work changes the
representation geometry (HypCD~\cite{liu2025hypcd}; BIA~\cite{bia2026}, which enlarges the spectrum of the batch
of class tokens), reduces bias (DebGCD~\cite{liu2025debgcd}, CausalGCD~\cite{causalgcd2026}), uses parts or
tokens (APL~\cite{apl2025}, TAR~\cite{tar2026}, SubViT~\cite{subvit2026}), aligns teacher and student attention
(FlipClass~\cite{flipclass2024}), or uses language (ATCG~\cite{atcg2026}).

\paragraph{Hierarchy in GCD.} SelEx~\cite{rastegar2024selex} uses hierarchical pseudo-labels from $k$-means;
SEAL~\cite{he2025seal} uses dataset hierarchies; HypCD relies on hyperbolic geometry without an explicit
hierarchy. FreqCone needs no labels, no text and no clustering to obtain its hierarchy.

\paragraph{Hyperbolic embeddings and entailment.} Hyperbolic image embeddings~\cite{khrulkov2020hyperbolic}
and clipped hyperbolic classifiers~\cite{guo2022clipped} made the Poincar\'e ball practical for vision.
Entailment cones~\cite{ganea2018cones} encode partial orders; MERU~\cite{desai2023meru},
HyCoCLIP~\cite{pal2025hycoclip} and HySAC~\cite{poppi2025hysac} use them between images and text, where the
text supplies the more general concept. GCD has no text; in FreqCone the blurred image plays that role.

\paragraph{Frequency.} FourEx~\cite{rastegar2026fourex} trains GCD with low- and high-pass views assigned to
separate latent slices; Feng and Ge~\cite{feng2025freqshift} study GCD under domain shift from a frequency
perspective. FreqCone differs in what the frequency views mean: they are ancestors in a partial order, placed
at depths, not extra positives or separate feature groups.

% =====================================================================================================
\section{Method}\label{sec:method}
\paragraph{Setup.} Let $\mathcal{D}=\mathcal{D}_l\cup\mathcal{D}_u$ with $K$ classes. A backbone
$\varphi$ and head $g$ give $z=g(\varphi(x))\in\mathbb{R}^n$. Following HypCD we use the Poincar\'e ball
$\mathbb{D}^n_c=\{p:\ c\|p\|^2<1\}$, feature clipping $\mathrm{C}_r(z)=\min(1,r/\|z\|)\,z$ and
\begin{equation}
u(x)=\exp_0^c\big(\mathrm{C}_r(g(\varphi(x)))\big),\qquad
\exp_0^c(v)=\tanh(\sqrt{c}\|v\|)\frac{v}{\sqrt{c}\|v\|}.
\end{equation}
With M\"obius addition $\oplus_c$, the distance and the depth are
\begin{equation}
d_c(p,q)=\tfrac{2}{\sqrt{c}}\artanh\big(\sqrt{c}\,\|(-p)\oplus_c q\|\big),\qquad
\delta_c(p)=d_c(0,p)=\tfrac{2}{\sqrt{c}}\artanh(\sqrt{c}\|p\|).
\end{equation}

\paragraph{Frequency views.} For a radial ideal low-pass filter with cut-off $f$ (cycles per image),
\begin{equation}
L_f(x)=\mathcal{F}^{-1}\big(\mathbf{1}[\|\omega\|\le f]\odot\mathcal{F}(G_\sigma * x)\big),
\end{equation}
where $G_\sigma$ is a small Gaussian blur against ringing. Given cut-offs $f_a<f_b$ we build two coarse
views $x^a=L_{f_a}(x)$ and $x^b=L_{f_b}(x)$ of the first augmented view $x$; $x'$ is the second view.

\paragraph{Depth by blur level.} Coarse views are placed at fixed depths and the network chooses only
their direction:
\begin{equation}
a=\frac{r_a}{\sqrt{c}}\frac{g(\varphi(x^a))}{\|g(\varphi(x^a))\|},\qquad
b=\frac{r_b}{\sqrt{c}}\frac{g(\varphi(x^b))}{\|g(\varphi(x^b))\|},\qquad r_a<r_b .
\end{equation}

\paragraph{Directional alignment.} Coarse views learn to point at the other sharp view, which receives
no gradient from this term ($\sg$ is the stop-gradient):
\begin{equation}
\Lb_{\mathrm{dir}}=\frac{1}{|B|}\sum_{i}\sum_{k\in\{a,b\}}
-\log\frac{\exp(\cos(k_i,\sg(u'_i))/\tau_k)}{\sum_j \exp(\cos(k_i,\sg(u'_j))/\tau_k)} .
\end{equation}

\paragraph{Frequency entailment.} With $\tilde p=\sqrt{c}\,p$, the cone half-aperture and the exterior angle
of Ganea et al.~\cite{ganea2018cones}, extended to curvature $-c$ by scaling, are
\begin{align}
\psi_c(p)&=\arcsin\Big(\min\Big\{1,\ K\frac{1-c\|p\|^2}{\sqrt{c}\|p\|}\Big\}\Big),\\
\Xi_c(p,q)&=\arccos\left(\frac{\langle\tilde p,\tilde q\rangle(1+\|\tilde p\|^2)-\|\tilde p\|^2(1+\|\tilde q\|^2)}
{\|\tilde p\|\,\|\tilde p-\tilde q\|\sqrt{1+\|\tilde p\|^2\|\tilde q\|^2-2\langle\tilde p,\tilde q\rangle}}\right),
\end{align}
and $E_c(p,q)=\max(0,\Xi_c(p,q)-\psi_c(p))$. The chain coarsest $\sqsubseteq$ coarse $\sqsubseteq$ sharp is
enforced by
\begin{equation}
\Lb_{\mathrm{ent}}=\frac{1}{|B|}\sum_i\big[E_c(a_i,b_i)+E_c(b_i,\bar u_i)\big],\qquad
\bar u_i=\sg(\|u_i\|)\frac{u_i}{\|u_i\|},
\end{equation}
so the cone can rotate the sharp view but never shrink its depth. With $K=0.1$, $r_a=0.12$ and $r_b=0.2$ the
half-apertures are about $55^\circ$ at $a$ and $29^\circ$ at $b$; both radii exceed the lower bound
$2K/(1+\sqrt{1+4K^2})\approx0.099$ of the cone's domain.

\paragraph{Scale contraction.}
\begin{equation}
\Lb_{\mathrm{con}}=\frac{1}{|B|(|B|-1)}\sum_{i\neq j}\Big([d_c(a_i,a_j)-d_c(b_i,b_j)]_+ + [d_c(b_i,b_j)-d_c(u_i,u_j)]_+\Big).
\end{equation}

\paragraph{Objective.} $\Lb=\Lb_{\mathrm{base}}(u,u')+\lambda_{\mathrm{dir}}\Lb_{\mathrm{dir}}
+\lambda_{\mathrm{ent}}\Lb_{\mathrm{ent}}+\lambda_{\mathrm{con}}\Lb_{\mathrm{con}}$, where
$\Lb_{\mathrm{base}}$ is the loss of the hyperbolic base method. Labels enter only through
$\Lb_{\mathrm{base}}$ on the sharp views. Inference uses the original image only. Algorithm~\ref{alg:train}
summarises one step.

\begin{algorithm}[t]
\caption{One FreqCone training step (added to a hyperbolic GCD base method).}\label{alg:train}
\begin{algorithmic}[1]
\Require batch $\{x_i\}$ with labels on the labelled subset; cut-offs $f_a<f_b$; radii $r_a<r_b$
\State $x,x'\gets$ two augmentations; $u,u'\gets\exp_0^c(\mathrm{C}_r(g(\varphi(\cdot))))$ \Comment{as the base method}
\State $x^a\gets L_{f_a}(x)$, $x^b\gets L_{f_b}(x)$
\State $a\gets \tfrac{r_a}{\sqrt c}\,\widehat{g(\varphi(x^a))}$, $b\gets \tfrac{r_b}{\sqrt c}\,\widehat{g(\varphi(x^b))}$ \Comment{$\hat v=v/\|v\|$}
\State $\Lb\gets\Lb_{\mathrm{base}}(u,u')+\lambda_{\mathrm{dir}}\Lb_{\mathrm{dir}}(a,b;\sg(u'))
+\lambda_{\mathrm{ent}}\Lb_{\mathrm{ent}}(a,b,\bar u)+\lambda_{\mathrm{con}}\Lb_{\mathrm{con}}(a,b,u)$
\State update $\varphi,g$ by back-propagating $\Lb$
\end{algorithmic}
\end{algorithm}

\paragraph{Cost.} Each image needs two extra forward and backward passes on coarse views, as in FourEx.
Because $L_f(x)$ has no energy above $f$, a coarse view can be resized to 112 pixels without aliasing whenever
$f$ is below the new Nyquist limit of 56 cycles per image; for a ViT-B/14 this quarters the number of tokens.
$\Lb_{\mathrm{con}}$ costs $O(|B|^2 n)$.

\paragraph{Attribution of novel samples.}\label{sec:attr} Let $\bar b_k$ be the mean coarse view of known
class $k$, projected back to radius $r_b/\sqrt c$. A sample $u$ is attributed to the family of
$\arg\min_k E_c(\bar b_k,u)$, with ties broken by the angle. This uses only the trained model and the labelled
set. We propose this rule for deployment but have not evaluated it: the mechanism study of
Section~\ref{sec:mech} measures attribution with the nearest labelled class prototype instead.

% =====================================================================================================
\section{Analysis}\label{sec:analysis}
\begin{proposition}[Low-pass filtration is a hierarchy]\label{prop:filtration}
Let $d_f(x,y)=\|L_f(x-y)\|_2$ (with $G_\sigma$ absorbed). Then (i) $L_fL_{f'}=L_{\min(f,f')}$;
(ii) $d_f(x,y)$ is non-decreasing in $f$; (iii) for every $\eta>0$, the connected components
$\mathcal{P}_f(\eta)$ of the graph linking images with $d_f\le\eta$ satisfy: if $f\le f'$, every block of
$\mathcal{P}_{f'}(\eta)$ is contained in a block of $\mathcal{P}_f(\eta)$.
\end{proposition}
\begin{proof}
(i) The masks are indicators of nested discs. (ii) By Parseval,
$d_f(x,y)^2=\frac{1}{HW}\sum_{\|\omega\|\le f}|\widehat{G_\sigma(x-y)}(\omega)|^2$, a sum of non-negative terms over
an index set that grows with $f$. (iii) A chain of steps with $d_{f'}\le\eta$ is, by (ii), also a chain with
$d_f\le\eta$.
\end{proof}

\begin{proposition}[Slicing after the exponential map entangles depth]\label{prop:slice}
Let $z=(z_A,z_B)$, $z_A\neq 0$, and $P_A$ the coordinate projection. Then
$\|P_A\exp_0^c(z)\|=\frac{\tanh(\sqrt{c}\|z\|)}{\sqrt{c}\|z\|}\|z_A\|\le\|\exp_0^c(z_A)\|$, with equality iff
$z_B=0$; the depth of the $A$-slice is strictly decreasing in $\|z_B\|$.
\end{proposition}
\begin{proof}
$h(t)=\tanh(t)/t$ is strictly decreasing on $(0,\infty)$ since $h'(t)=(t\,\mathrm{sech}^2t-\tanh t)/t^2<0$,
equivalently $\sinh 2t>2t$; and $\|z\|\ge\|z_A\|$. Depth is increasing in the Euclidean norm.
\end{proof}
Hence, when frequency bands are assigned to latent slices after a full-vector exponential map, the depth of one
band depends on the energy of the other, so depth cannot encode the frequency level.

\begin{lemma}[Cones imply depth ordering]\label{lem:depth}
If $E_c(p,q)=0$ and $\psi_c(p)\le\pi/2$, then $\delta_c(q)\ge\delta_c(p)$, strictly if $\psi_c(p)<\pi/2$ and
$q\neq p$.
\end{lemma}
\begin{proof}
Let $\gamma$ be the geodesic from $p$ to $q$ and $f(t)=\delta_c(\gamma(t))$. Hyperbolic space is
CAT(0)~\cite{bridson1999metric}, so $f$ is convex. The gradient of $\delta_c$ at $p$ is the outward radial unit
vector and the model is conformal, hence $f'(0)=\cos\Xi_c(p,q)$, where $\Xi_c(p,q)=\pi-\angle Opq$. Zero energy
gives $\Xi_c\le\psi_c\le\pi/2$, so $f'(0)\ge0$ and the convex $f$ is non-decreasing.
\end{proof}

\begin{proposition}[Contraction transfers the hierarchy]\label{prop:contraction}
If $\Lb_{\mathrm{con}}=0$ on a set of images, then for every $\eta$ the $\eta$-components of $\{u_i\}$
refine those of $\{b_i\}$, which refine those of $\{a_i\}$.
\end{proposition}
\begin{proof}Identical to Proposition~\ref{prop:filtration}(iii) with $d_c$ in place of $d_f$.\end{proof}

\begin{proposition}[Cones bound angular deviation]\label{prop:angle}
If $E_c(p,q)=0$ then the angle at the origin satisfies $\angle(p,q)<\psi_c(p)$.
Consequently, for images $i,j$ with $\angle(a_i,a_j)\le\eta$ whose sharp views lie in their coarse cones,
$\angle(u_i,u_j)<\psi_c(a_i)+\eta+\psi_c(a_j)$.
\end{proposition}
\begin{proof}
In the hyperbolic triangle $O,p,q$ the angle at $p$ is $\pi-\Xi_c(p,q)\ge\pi-\psi_c(p)$; the angle sum is below
$\pi$, so the angle at $O$ is below $\psi_c(p)$. Geodesics through $O$ are straight lines, so this is the
Euclidean angle. The second statement is the triangle inequality for angles on the sphere.
\end{proof}

\paragraph{Remark (when a cone cannot act).} If $\delta_c(q)<\delta_c(p)$, Lemma~\ref{lem:depth} gives
$E_c(p,q)>0$ for every direction of $q$. Moreover, for $q$ on the ray of $p$ the geodesic from $p$ to $q$ points
at the origin, so $\Xi_c(p,q)=\pi$, the largest possible value. Hence, if the norm of $q$ is held fixed, any
descent on $E_c$ rotates $q$ away from the direction of $p$. For the depths met in Section~\ref{sec:mech}
($\delta_c(b)=1.81$, $\delta_c(u)=1.3$) the minimising angle is $42.7^\circ$ and the minimum of $\Xi_c$ is
$135^\circ$, far above $\psi_c(b)=28.7^\circ$; for $\delta_c(u)=4.6$ the minimising angle is $0$.

\paragraph{Numerical checks.} \texttt{verify\_math.py} checks the closed form of $\Xi_c$ against the angle
computed from the geodesic (error $7\times10^{-12}$), the curvature scaling, Lemma~\ref{lem:depth} on random
pairs inside cones (minimum depth gain $+0.022$), Proposition~\ref{prop:angle} (on every sampled pair the angle at the
origin is at least $0.020$\,rad below $\psi_c$), Proposition~\ref{prop:slice} (slice norms shrink by 19\% on average
for random vectors) and Proposition~\ref{prop:filtration} on random images.

\paragraph{What the analysis does not show.} Proposition~\ref{prop:filtration} is about pixel distances; whether
the frequency hierarchy matches the semantic one is an empirical question that depends on the dataset.
Proposition~\ref{prop:angle} bounds how far siblings can drift apart; it says nothing about how close they are
pushed, which is where fine-grained accuracy can be lost (Section~\ref{sec:mech}).

% =====================================================================================================
\section{Mechanism study}\label{sec:mech}
This study runs on CPU with synthetic data. It tests whether the mechanism behaves as the analysis predicts
and what it costs; it does not replace the SSB evaluation of Section~\ref{sec:ssb}.

\paragraph{Data (SpecHier-48).} Eight parents times six children give 48 classes of $16\times16$ images. A
parent is a smooth layout of three Gaussian blobs (low frequency); a child adds an oriented grating at about
five cycles per image (high frequency). The six orientations are shared by all parents and permuted per parent,
so a class is identified only by the combination. Images are shifted by up to 3 pixels, the grating phase,
contrast and brightness are random, and white noise with $\sigma=0.35$ is added. Logistic regression on raw
pixels reaches 14.4\% on the 48 classes and 83.3\% on the parents.
Each parent has three known and three novel children (24/24 classes, 60 images per class); half of the
known-class images are labelled (720), and the remaining 2160 images are unlabelled. Twenty further images per
known class form a validation set that is never trained on.

\paragraph{Model and training.} A four-layer CNN (widths 32-64-96-128) with a $128\to32$ head is trained from
scratch for 40 epochs with AdamW (learning rate $2\cdot10^{-3}$, batch 192). Hyperbolic methods use $c=0.05$,
$r=2.3$ and the hybrid angle-to-distance loss of HypCD; FreqCone uses $f_a=1.5$, $f_b=3$ cycles per image and the
other hyper-parameters of Section~\ref{sec:method} ($\lambda_{\mathrm{dir}}=0.5$, $\lambda_{\mathrm{con}}=0.05$,
$\tau_a=0.5$, $\tau_b=0.3$).

\paragraph{Methods.} GCD (spherical supervised plus self-supervised contrastive loss); HypCD; a FourEx-style
partition on HypCD (one half of the latent vector trained with low-pass views, the other with high-pass views);
low-pass positives on HypCD (the coarse views are ordinary positives of the second view); and FreqCone on HypCD.

\paragraph{Metrics.} Semi-supervised $k$-means on backbone features, scored with one Hungarian matching over all
unlabelled images (All/Old/New). Two structural metrics: parent accuracy ($k$-means with 8 clusters) and
novel-to-parent attribution, the share of novel images whose nearest labelled class prototype has the same
parent (chance 12.5\%).

\paragraph{Protocol.} Seed 0 is the development seed; seeds 1--3 are held out, and each seed draws new data and
a new initialisation. On seed 0 we sweep $\lambda_{\mathrm{ent}}\in\{0,0.03,0.1,0.3,1\}$ and select
$\lambda^\star$ by the clustering accuracy on the known-class validation images ($k$-means with one cluster per
known class, Hungarian matching), breaking ties toward the smaller weight. The rule was fixed before the sweep
was run; no novel-class label is used for selection. A first criterion, nearest-prototype accuracy on the
validation images, reached 100\% for $\lambda_{\mathrm{ent}}=0$ and could not rank the weights, so it was
replaced before any other weight was evaluated. On the held-out seeds we also report $\lambda_{\mathrm{ent}}=0$
(FreqCone without $\Lb_{\mathrm{ent}}$, an ablation) and $\lambda_{\mathrm{ent}}=1$ (the strong end of the
trade-off); neither was used to choose anything.

\subsection{Design log on the development seed}
Table~\ref{tab:designlog} lists every FreqCone variant we trained with $\lambda_{\mathrm{ent}}=1$. All three
produce strong parent structure, but the first two damage fine-grained accuracy for identifiable reasons.
In v0 the sharp views were pulled toward the origin (depth \logDepthVzero{} against the clipping limit
4.60), where cones approach half-spaces. Detaching the sharp-view norm inside $\Lb_{\mathrm{ent}}$ (v1) did not
help (depth \logDepthVone), so the inward pull did not come from $\Lb_{\mathrm{ent}}$. The causes were
$\Lb_{\mathrm{dir}}$, which pulled the sharp view toward a coarse view that only knows the parent and so merged
siblings, and the shared encoder, which shrank every norm to widen the cones. v2 fixes the coarse depths and
stops the gradient to the sharp view in $\Lb_{\mathrm{dir}}$; sharp views return to the boundary
(depth \logDepthVtwo). With $\lambda_{\mathrm{ent}}=1$, however, v2 still loses fine-grained accuracy, which
motivates selecting the weight.

\begin{table}[t]
\centering\small
\IfFileExists{tab_designlog.tex}{\resizebox{\linewidth}{!}{% generated by make_tables.py from pilot/*.json (seed 0)
\begin{tabular}{llcccc}
\toprule
Variant & Change & All / Old / New & Parent & Novel$\to$parent & depth$(u)$ \\
\midrule
HypCD (base) & -- & 73.6 / 68.2 / 76.3 & 34.6 & 54.0 & 4.60 \\
FreqCone v0 & initial formulation & 56.7 / 55.4 / 57.3 & 78.8 & 88.6 & 1.25 \\
FreqCone v1 & stop-gradient on sharp-view norm in $\mathcal{L}_{\mathrm{ent}}$ & 29.9 / 25.4 / 32.1 & 97.7 & 96.5 & 0.89 \\
FreqCone v2 & fixed coarse depths; stop-gradient on $u'$ in $\mathcal{L}_{\mathrm{dir}}$ & 48.9 / 42.8 / 51.9 & 100.0 & 97.0 & 4.60 \\
\bottomrule
\end{tabular}
}}{}
\caption{Development seed 0, $\lambda_{\mathrm{ent}}=1$. Depth is the mean $\delta_c(u)$ of sharp views; the
clipping radius limits it to 4.60. For reference, spherical GCD reaches 80.0 / 73.6 / 83.2 with 22.8\% parent
accuracy and 23.5\% novel-to-parent attribution.}\label{tab:designlog}
\end{table}

\subsection{Selecting the entailment weight}
On the development seed (Figure~\ref{fig:sweep}) the rule selects $\lambda^\star=\selLambda$, whose validation
clustering accuracy (\devValcb\%) is the highest of the five weights. Validation and test disagree at
$\lambda_{\mathrm{ent}}=0.1$: test All accuracy is higher (\devAllc{} against \devAllb) and parent accuracy rises
from \devParb{} to \devParc\%, while validation accuracy falls to \devValcc\%. A plausible reason is that the
validation set contains only known classes, so it cannot reward a parent structure that helps to group novel
images. We keep $\lambda^\star$ as selected and report $\lambda_{\mathrm{ent}}=0.1$ only in the sensitivity
analysis. From $\lambda_{\mathrm{ent}}=0.3$ on, both structural metrics keep rising while All accuracy falls, to
\devAlle\% at $\lambda_{\mathrm{ent}}=1$. This matches the limitation stated after
Proposition~\ref{prop:angle}: cones bound how far siblings drift apart, but nothing keeps them apart.

\begin{figure}[t]
\centering
\IfFileExists{lambda_sweep.dat}{%
\begin{tikzpicture}
\begin{groupplot}[group style={group size=2 by 1, horizontal sep=1.6cm},
  width=0.47\linewidth, height=5.2cm, xtick={0,1,2,3,4}, xticklabels={0,0.03,0.1,0.3,1},
  xlabel={$\lambda_{\mathrm{ent}}$}, ymin=0, ymax=100, grid=major, grid style={gray!20},
  legend style={font=\scriptsize, draw=none, fill=white, fill opacity=0.8, text opacity=1}, tick label style={font=\scriptsize}]
\nextgroupplot[title={\small (a) fine-grained clustering}, ylabel={accuracy (\%)}, legend pos=south west]
\addplot[blue!70!black, thick, mark=*] table[x=idx, y=all] {lambda_sweep.dat};
\addplot[gray, thick, dashed, mark=square*] table[x=idx, y=valc] {lambda_sweep.dat};
\addplot[gray, dotted, thick, domain=0:4, samples=2] {\logAllHyp};
\legend{All (test), known-class validation, HypCD All}
\nextgroupplot[title={\small (b) hierarchical consistency}, legend pos=south east]
\addplot[blue!70!black, thick, mark=*] table[x=idx, y=n2p] {lambda_sweep.dat};
\addplot[gray, thick, dashed, mark=square*] table[x=idx, y=parent] {lambda_sweep.dat};
\addplot[gray, dotted, thick, domain=0:4, samples=2] {\logNtpHyp};
\legend{novel-to-parent, parent accuracy, HypCD novel-to-parent}
\end{groupplot}
\end{tikzpicture}}{\fbox{lambda sweep not available}}
\caption{Development seed 0. Raising $\lambda_{\mathrm{ent}}$ trades fine-grained accuracy (a) for hierarchical
consistency (b). The validation curve in (a) uses only known-class images and is the selection criterion.}
\label{fig:sweep}
\end{figure}

\subsection{Held-out seeds}
\begin{table}[t]
\centering\small
\IfFileExists{tab_heldout.tex}{% generated by make_tables.py from results/*.json
\begin{tabular}{lccccc}
\toprule
Method & All & Old & New & Parent & Novel$\to$parent \\
\midrule
GCD (spherical) & 76.0{\scriptsize$\pm$11.5} & \textbf{78.7{\scriptsize$\pm$13.7}} & 74.7{\scriptsize$\pm$11.4} & 24.6{\scriptsize$\pm$2.1} & 22.2{\scriptsize$\pm$4.6} \\
HypCD & 75.4{\scriptsize$\pm$9.6} & 73.1{\scriptsize$\pm$14.4} & 76.6{\scriptsize$\pm$7.3} & 30.8{\scriptsize$\pm$12.4} & 32.2{\scriptsize$\pm$26.7} \\
FourEx-style partition & 70.9{\scriptsize$\pm$2.6} & 73.7{\scriptsize$\pm$3.3} & 69.5{\scriptsize$\pm$2.5} & 26.0{\scriptsize$\pm$6.2} & 28.7{\scriptsize$\pm$22.2} \\
Low-pass positives & 44.5{\scriptsize$\pm$20.2} & 45.9{\scriptsize$\pm$25.4} & 43.8{\scriptsize$\pm$17.7} & \textbf{74.7{\scriptsize$\pm$31.8}} & \textbf{85.1{\scriptsize$\pm$20.1}} \\
FreqCone without $\mathcal{L}_{\mathrm{ent}}$ ($\lambda_{\mathrm{ent}}{=}0$) & 75.8{\scriptsize$\pm$12.0} & 72.4{\scriptsize$\pm$11.1} & \textbf{77.5{\scriptsize$\pm$14.0}} & 30.1{\scriptsize$\pm$13.0} & 28.6{\scriptsize$\pm$31.8} \\
FreqCone ($\lambda^\star$) & \textbf{76.5{\scriptsize$\pm$11.6}} & 75.8{\scriptsize$\pm$6.7} & 76.8{\scriptsize$\pm$14.1} & 29.9{\scriptsize$\pm$13.1} & 29.4{\scriptsize$\pm$33.0} \\
FreqCone ($\lambda_{\mathrm{ent}}{=}1$) & 69.8{\scriptsize$\pm$26.2} & 67.2{\scriptsize$\pm$21.3} & 71.2{\scriptsize$\pm$28.7} & 44.5{\scriptsize$\pm$38.2} & 41.1{\scriptsize$\pm$49.8} \\
\bottomrule
\end{tabular}
}{}
\caption{Held-out seeds 1--3, mean {\scriptsize$\pm$} standard deviation. Bold marks the best mean per column.}
\label{tab:heldout}
\end{table}
\begin{table}[t]
\centering\small
\IfFileExists{tab_paired.tex}{% generated by make_tables.py
\begin{tabular}{lcccc}
\toprule
vs.\ HypCD (same seed) & $\Delta$All & $\Delta$New & $\Delta$Parent & $\Delta$Novel$\to$parent \\
\midrule
GCD (spherical) & +0.6 (1/3) & -1.9 (1/3) & -6.2 (2/3) & -10.0 (2/3) \\
FourEx-style partition & -4.5 (1/3) & -7.0 (0/3) & -4.8 (0/3) & -3.6 (1/3) \\
Low-pass positives & -30.9 (0/3) & -32.8 (0/3) & +43.9 (3/3) & +52.9 (3/3) \\
FreqCone without $\mathcal{L}_{\mathrm{ent}}$ ($\lambda_{\mathrm{ent}}{=}0$) & +0.4 (1/3) & +0.9 (1/3) & -0.7 (2/3) & -3.6 (1/3) \\
FreqCone ($\lambda^\star$) & +1.1 (2/3) & +0.3 (2/3) & -0.9 (1/3) & -2.8 (1/3) \\
FreqCone ($\lambda_{\mathrm{ent}}{=}1$) & -5.6 (2/3) & -5.4 (2/3) & +13.7 (2/3) & +8.9 (1/3) \\
\bottomrule
\end{tabular}
}{}
\caption{Paired difference to HypCD on the same seed, in points; in brackets, the number of seeds on which the
method is better.}\label{tab:paired}
\end{table}
Table~\ref{tab:heldout} reports the held-out seeds and Table~\ref{tab:paired} the paired differences to
HypCD. At $\lambda^\star$, FreqCone exceeds HypCD by \dAllSel{} points of All accuracy and wins on \wAllSel{}
seeds, while the standard deviation across seeds is about 10 points; against FreqCone without
$\Lb_{\mathrm{ent}}$ the difference is \pairSelVsWzeroAll{} (\pairwSelVsWzeroAll{} seeds). At the selected weight
this study therefore shows no reliable gain in clustering accuracy, and novel-to-parent attribution does not
improve either (\dNtpSel{} points).

Two comparisons are clearer. Treating blurred views as ordinary positives (H4) costs \dAllNaive{} points of All
accuracy and loses on every seed, while raising novel-to-parent attribution by \dNtpNaive{} points; FreqCone with
$\lambda_{\mathrm{ent}}=1$ is \pairWoneVsNaiveAll{} points more accurate on all three seeds but
\pairWoneVsNaiveNtp{} points lower in attribution, so the two sit at different points of the same trade-off.
Assigning frequency bands to latent slices on the same hyperbolic base (H3) costs \dAllFour{} points of All
and \dNewFour{} points of New accuracy relative to HypCD; FreqCone at $\lambda^\star$ is \pairSelVsFourAll{}
points above it (\pairwSelVsFourAll{} seeds).

\paragraph{Two depth regimes.} The large spread across seeds has a specific cause (Table~\ref{tab:regimes}).
On seeds 0 and 3, HypCD pushes sharp views to the clipping boundary (depth 4.60), deeper than the coarse view
$b$ (fixed at depth 1.81); cones are satisfiable and the share of sharp views inside their cone grows with
$\lambda_{\mathrm{ent}}$, which produces the trade-off of Figure~\ref{fig:sweep}. On seeds 1 and 2, HypCD keeps
sharp views at depth about 1.2, shallower than $b$. By Lemma~\ref{lem:depth} no sharp view can then lie in its
cone (0\% at every weight), and because v2 detaches the sharp-view norm inside $\Lb_{\mathrm{ent}}$, the cone can
only rotate the sharp view; by the Remark in Section~\ref{sec:analysis} it rotates it away from its own coarse
view. The measured angle between $b$ and $u$ on these seeds is about $45^\circ$, against the predicted
$42.7^\circ$. In this regime $\Lb_{\mathrm{ent}}$ does the opposite of its purpose: accuracy can rise (seed 1:
$+10.4$ points at $\lambda_{\mathrm{ent}}=1$) but no hierarchy forms. Averages over seeds therefore mix two
different mechanisms. The direct fix is to let the cone move the sharp-view norm (v3, next subsection): with
coarse depths fixed, moving a sharp view outward along its ray never increases $\Xi_c$ (checked on
$2\times10^5$ random configurations), so the fix should not reopen the collapse toward the origin.

\begin{table}[t]
\centering\small
\IfFileExists{tab_regimes.tex}{\resizebox{\linewidth}{!}{% generated by make_tables.py: seed 0 = development, others held out
\begin{tabular}{lc|cccc|cccc}
\toprule
 & HypCD & \multicolumn{4}{c|}{FreqCone v2, $\lambda^\star$} & \multicolumn{4}{c}{FreqCone v2, $\lambda_{\mathrm{ent}}{=}1$} \\
Seed & $\delta_c(u)$ & $\delta_c(u)$ & in cone & $\Delta$All & $\Delta$N$\to$P & $\delta_c(u)$ & in cone & $\Delta$All & $\Delta$N$\to$P \\
\midrule
0 (dev) & 4.60 & 4.60 & 7\% & +0.7 & +12.8 & 4.60 & 82\% & -24.7 & +43.0 \\
1 & 1.17 & 1.32 & 0\% & +5.1 & -2.2 & 1.32 & 0\% & +10.4 & -1.8 \\
2 & 1.15 & 1.25 & 0\% & +1.1 & -11.9 & 1.30 & 0\% & +0.3 & -8.1 \\
3 & 4.60 & 4.60 & 3\% & -3.0 & +5.6 & 4.60 & 78\% & -27.3 & +36.5 \\
\bottomrule
\end{tabular}
}}{}
\caption{Per-seed diagnostics of FreqCone v2. $\delta_c(u)$ is the mean depth of sharp views (the coarse view
$b$ sits at 1.81); ``in cone'' is the share of sharp views inside the cone of their coarse view $b$; $\Delta$ are
differences to HypCD on the same seed, in points.}\label{tab:regimes}
\end{table}

\subsection{v3 on fresh seeds}
v3 removes the stop-gradient on the sharp-view norm in $\Lb_{\mathrm{ent}}$ and changes nothing else. Because
seeds 1--3 motivated it, we evaluate it on three fresh seeds (4--6), with the weight selected for v2
($\lambda^\star=\selLambda$) and with $\lambda_{\mathrm{ent}}=1$; no weight is re-selected. Two questions were
written into the run script before any v3 result existed (Table~\ref{tab:v3}).

\emph{Q1: does v3 put sharp views inside their cones where v2 cannot?} Only at $\lambda_{\mathrm{ent}}=1$. At
$\lambda^\star$ the sharp views of v3 stay at depth about 1.3 and none lies inside its cone, so the outward force
of $\Lb_{\mathrm{ent}}$ is too weak against the rest of the objective. At $\lambda_{\mathrm{ent}}=1$, v3 moves them
to the boundary on every fresh seed, \VthreeConeRange{} of them lie inside their cone, and novel-to-parent
attribution reaches 93.8--99.5\%. Seed 4 also shows that the shallow regime is not inherited from the base model
alone: HypCD reaches the boundary there, while v2 and v3 at $\lambda^\star$ stay shallow.

\emph{Q2: accuracy.} On these seeds v2 at $\lambda^\star$ is \vdAllVtwo{} points of All accuracy above HypCD
(\vwAllVtwo{} seeds) but \vdNtpVtwo{} points lower in novel-to-parent attribution (\vwNtpVtwo{}), as expected from a
cone that pushes samples away from their parents. At $\lambda^\star$, v3 is not better than v2 (\vdAllVthreeVsVtwo{} points of All accuracy,
better on \vwAllVthreeVsVtwo{} seeds) and differs from HypCD by \vdAllVthree{} (\vwAllVthree{}). With an active
cone ($\lambda_{\mathrm{ent}}=1$) it differs from HypCD by \vdAllVthreeWone{} points of All accuracy and
\vdNtpVthreeWone{} points of novel-to-parent attribution. The fix repairs the mechanism but not the outcome:
Proposition~\ref{prop:angle} caps the angle between siblings that share a coarse direction at about
$2\psi_c(b)\approx57^\circ$, and once the cone is active this cap binds and fine-grained accuracy falls. Wider
cones (a larger $K$) relax the cap; we have not tested them.

\begin{table}[t]
\centering\small
\IfFileExists{tab_v3.tex}{% generated by make_tables.py from results_v3/*.json (fresh seeds 4-6)
\begin{tabular}{llccccc}
\toprule
Seed & Method & $\delta_c(u)$ & in cone & All & New & Novel$\to$parent \\
\midrule
4 & HypCD & 4.60 & -- & 71.2 & 70.8 & 78.1 \\
 & v2, $\lambda^\star$ & 1.25 & 0\% & 74.7 & 73.7 & 25.5 \\
 & v3, $\lambda^\star$ & 1.30 & 0\% & 73.3 & 74.9 & 25.9 \\
 & v3, $\lambda_{\mathrm{ent}}{=}1$ & 4.60 & 83\% & 43.9 & 48.1 & 99.5 \\
\midrule
5 & HypCD & 1.29 & -- & 71.4 & 71.8 & 28.3 \\
 & v2, $\lambda^\star$ & 1.28 & 0\% & 76.4 & 78.2 & 23.5 \\
 & v3, $\lambda^\star$ & 1.33 & 0\% & 70.7 & 70.4 & 26.8 \\
 & v3, $\lambda_{\mathrm{ent}}{=}1$ & 4.60 & 80\% & 37.0 & 38.8 & 99.1 \\
\midrule
6 & HypCD & 1.16 & -- & 84.6 & 88.2 & 12.8 \\
 & v2, $\lambda^\star$ & 1.28 & 0\% & 83.5 & 87.4 & 10.7 \\
 & v3, $\lambda^\star$ & 1.32 & 0\% & 84.2 & 81.8 & 12.2 \\
 & v3, $\lambda_{\mathrm{ent}}{=}1$ & 4.60 & 75\% & 52.0 & 51.0 & 93.8 \\
\bottomrule
\end{tabular}
}{}
\caption{FreqCone v3 on fresh seeds 4--6. v3 differs from v2 only in letting $\Lb_{\mathrm{ent}}$ move the norm of
the sharp view. ``In cone'' is the share of sharp views inside the cone of their coarse view $b$.}\label{tab:v3}
\end{table}

\subsection{What the study shows and what it does not}
\paragraph{What the study supports.} The mechanism behaves as the analysis predicts. When sharp views are deeper
than their coarse views, the cone is satisfiable and produces a strong hierarchy, with novel-to-parent
attribution up to 97--99\%. When they are shallower, no sample can enter its cone (Lemma~\ref{lem:depth}) and a
norm-detached cone pushes samples away from their parents (Remark in Section~\ref{sec:analysis}). Entailment is
far less destructive than treating blurred views as positives, and encoding frequency as depth does better than a
latent partition on the same hyperbolic base.

\paragraph{What it does not support.} At the weight selected on validation data, FreqCone does not improve
clustering accuracy reliably over HypCD, and whenever the cone is active the hierarchy costs fine-grained
accuracy. In this setting FreqCone is a controllable trade-off between hierarchical consistency and fine-grained
accuracy, not an accuracy improvement.

\paragraph{Caveats.} The data are synthetic and their frequency hierarchy is the semantic one by construction; the
encoder is a small CNN trained from scratch, not a pre-trained ViT; each comparison uses three seeds and the
variance across seeds is large; novel-to-parent attribution is measured with class prototypes, not with the cone
rule of Section~\ref{sec:attr}.

\paragraph{Implications for the SSB evaluation.} Log the sharp-view depth and the share of samples inside their
cone during training; sweep the cone aperture $K$ together with $\lambda_{\mathrm{ent}}$; and report family-level
attribution next to All, Old and New accuracy.

% =====================================================================================================
\section{Protocol for the SSB evaluation}\label{sec:ssb}
No number in this section has been run. Backbone: DINOv2 ViT-B/14~\cite{oquab2024dinov2}. Base methods:
Hyp-SimGCD and Hyp-SelEx from the official HypCD code, re-run in the same environment, each with and without
FreqCone; FourEx from its official code for the latent-partition comparison. Data: CUB, Stanford Cars and
FGVC-Aircraft from SSB, with Herbarium19, CIFAR-100 and ImageNet-100 to check that FreqCone does no harm on
coarse data. Scoring: one Hungarian matching over all unlabelled images, last epoch, three seeds, mean and
standard deviation. Hyper-parameters, including $\lambda_{\mathrm{ent}}$, are selected on known-class validation
data as in Section~\ref{sec:mech}.

\paragraph{Decision criteria, fixed in advance.} FreqCone is called effective if its mean New accuracy exceeds
that of the base method by at least 1.0 point on at least two of the three SSB datasets, for both base methods,
with a paired win on at least two of three seeds, and if All accuracy drops by no more than 0.5 points on any
dataset. Negative results will be reported, in particular on Stanford Cars, where FourEx reports lower New
accuracy than SelEx (73.2 against 76.7).

\paragraph{Ablations.} Removing $\Lb_{\mathrm{dir}}$, $\Lb_{\mathrm{ent}}$ or $\Lb_{\mathrm{con}}$; entailment
against low-pass positives; depth encoding against a latent partition on the same hyperbolic base; one, two or
three blur levels; cone aperture $K$; Euclidean against hyperbolic cones; 112-pixel coarse views.

\paragraph{Analyses.} Mean depth per blur level and the share of ordered triples
$\delta_c(a)\le\delta_c(b)\le\delta_c(u)$; whether the frequency hierarchy matches the semantic one, measured
with the manufacturer and family labels of FGVC-Aircraft and the makes of Stanford Cars; and novel-to-family
attribution with the rule of Section~\ref{sec:attr}.

% =====================================================================================================
\section{Limitations}
The evidence for FreqCone in this draft comes from proofs and from a synthetic mechanism study; whether it
improves GCD on natural images is decided by the protocol of Section~\ref{sec:ssb}, which has not been run.
SpecHier-48 is built so that the frequency hierarchy is the semantic one; on real data, blurred views can be
dominated by background or pose rather than by family. Cones constrain directions but not how close siblings
are, so a large entailment weight can merge them, as Figure~\ref{fig:sweep} shows. The coarse views roughly
double the training cost unless they are processed at reduced resolution.

% =====================================================================================================
\begin{thebibliography}{99}\small
\bibitem{vaze2022gcd} S.~Vaze, K.~Han, A.~Vedaldi, A.~Zisserman. Generalized Category Discovery. CVPR, 2022.
\bibitem{vaze2022ssb} S.~Vaze, K.~Han, A.~Vedaldi, A.~Zisserman. Open-Set Recognition: a Good Closed-Set Classifier is All You Need? ICLR, 2022.
\bibitem{wen2023simgcd} X.~Wen, B.~Zhao, X.~Qi. Parametric Classification for Generalized Category Discovery: A Baseline Study. ICCV, 2023.
\bibitem{vaze2023mugcd} S.~Vaze, A.~Vedaldi, A.~Zisserman. No Representation Rules Them All in Category Discovery. NeurIPS, 2023.
\bibitem{rastegar2024selex} S.~Rastegar, M.~Salehi, Y.~M.~Asano, H.~Doughty, C.~G.~M.~Snoek. SelEx: Self-Expertise in Fine-Grained Generalized Category Discovery. ECCV, 2024.
\bibitem{liu2025hypcd} Y.~Liu, Z.~He, K.~Han. Hyperbolic Category Discovery. CVPR, 2025.
\bibitem{liu2025debgcd} Y.~Liu, K.~Han. DebGCD: Debiased Learning with Distribution Guidance for Generalized Category Discovery. ICLR, 2025.
\bibitem{rastegar2026fourex} S.~Rastegar, M.~Ghadimi~Atigh, P.~Mettes, Y.~M.~Asano, C.~G.~M.~Snoek. Fourier Self-Supervision for Fine-Grained Generalized Category Discovery. arXiv:2608.08963, 2026.
\bibitem{he2025seal} Z.~He, Y.~Liu, K.~Han. SEAL: Semantic-Aware Hierarchical Learning for Generalized Category Discovery. NeurIPS, 2025.
\bibitem{bia2026} Tang et~al. Bures-Isotropy Alignment: Manifold Learning of Generalized Category Discovery. ICLR, 2026.
\bibitem{flipclass2024} H.~Lin et~al. Flipped Classroom: Aligning Teacher Attention with Student in Generalized Category Discovery. NeurIPS, 2024.
\bibitem{apl2025} Dai et~al. Adaptive Part Learning for Fine-Grained Generalized Category Discovery. CVPR, 2025.
\bibitem{tar2026} Yang et~al. TAR: Token-Aware Refinement for Fine-grained Generalized Category Discovery. CVPR, 2026.
\bibitem{atcg2026} Han et~al. Learning Like Humans: Analogical Concept Learning for Generalized Category Discovery. CVPR, 2026.
\bibitem{causalgcd2026} Feng et~al. Seeing Through the Shift: Causality-Inspired Robust Generalized Category Discovery. CVPR, 2026.
\bibitem{subvit2026} Zhu et~al. Subtoken Vision Transformer for Fine-grained Recognition. arXiv, 2026.
\bibitem{feng2025freqshift} Feng, Ge. Generalized Category Discovery under Domain Shift: A Frequency Domain Perspective. NeurIPS, 2025.
\bibitem{ganea2018cones} O.-E.~Ganea, G.~B\'ecigneul, T.~Hofmann. Hyperbolic Entailment Cones for Learning Hierarchical Embeddings. ICML, 2018.
\bibitem{khrulkov2020hyperbolic} V.~Khrulkov, L.~Mirvakhabova, E.~Ustinova, I.~Oseledets, V.~Lempitsky. Hyperbolic Image Embeddings. CVPR, 2020.
\bibitem{guo2022clipped} Y.~Guo, X.~Wang, Y.~Chen, S.~X.~Yu. Clipped Hyperbolic Classifiers Are Super-Hyperbolic Classifiers. CVPR, 2022.
\bibitem{desai2023meru} K.~Desai, M.~Nickel, T.~Rajpurohit, J.~Johnson, R.~Vedantam. Hyperbolic Image-Text Representations. ICML, 2023.
\bibitem{pal2025hycoclip} A.~Pal et~al. Compositional Entailment Learning for Hyperbolic Vision-Language Models. ICLR, 2025.
\bibitem{poppi2025hysac} T.~Poppi et~al. Hyperbolic Safety-Aware Vision-Language Models. CVPR, 2025.
\bibitem{oquab2024dinov2} M.~Oquab et~al. DINOv2: Learning Robust Visual Features without Supervision. TMLR, 2024.
\bibitem{bridson1999metric} M.~R.~Bridson, A.~Haefliger. Metric Spaces of Non-Positive Curvature. Springer, 1999.
\end{thebibliography}
\end{document}
